\subsection{Caracterización del decodificador}

Para cada peso entre 0 y 4, se recorre la totalidad de los patrones
de error correspondientes a dicho peso. Estos vectores son generados
en el golden model \href{https://github.com/AbelCorvalan0/fec-lab2/blob/main/model/golay_decoder.ipynb}{\textbf{golay decoder.ipynb}}.

Luego, estos vectores son procesados por un objeto decodificador para
obtener el archivo de vectores \href{https://github.com/AbelCorvalan0/fec-lab2/blob/main/model/outputs/decoder/decoder_top_testing/codeword_with_errors.txt}{\textbf{codeword with errors.txt}}.

Este archivo tiene el siguiente formato:
\begin{table}[h!]
    \centering
    \begin{tabular}{|c|c|c|c|c|}
        \hline
        1. Vector recibido &
        2. Mensaje &
        3. Máscara de error &
        4. Palabra corregida &
        5. Palabra incorregible \\
        \hline
    \end{tabular}
    %\caption{Descripción de las columnas de salida del decodificador.}
    \label{tab:decoder_output}
\end{table}

Con estos vectores se valido el modelo del decoder obteniendo exitosamente la tabla esperada de casos segun el peso de error:

\begin{figure}[H]
    \centering
    \includegraphics[scale=0.4]{decoder-error-table.png}
    \hfill
    \includegraphics[scale=0.4]{decoder-error-model.png}
    \caption{Tabla de casos de decodificacion segun peso de error (modelo).}
    \label{fig:decoder-error-table}
\end{figure}

Además se genera un conjunto de vectores por cada submódulo, todos los archivos se encuentran en \href{https://github.com/AbelCorvalan0/fec-lab2/tree/main/model/outputs/decoder}{\textbf{(vectores)}}.